% Korean Noether paper 14 — fresh producer translation, unchecked.
% Initial translation authority: ED0005 lines 9105--9283.
% Current scope-rebase authority: ED0011 / NOETH-DE-AUTH-v047-20260808.
% Authority SHA-256: 1A44F967B29972E8F99E5C323A479162AD82A23FC457395915A4BB9DDF51AD41.
% Normalized source span: 63,758 bytes; SHA-256 9DB240C8CC76AFBE3C95853AAB0F58BA9965EE5AFC301B2EF52F072033854964.
% Translation only; independent source, Korean-language, render, and certification checks remain external.

\section*{14. 한 변수 대수함수의 산술이론과 다른 이론들 및 수체론과의 관계}

\begin{center}
괴팅겐의 에미 뇌터 보고\\[0.6em]
\emph{Jahresber. d. D. Math.-Verein. 28 (1919), 182--203쪽}
\end{center}

이 보고서는 대수함수에 관한 개별 이론들을 서로 이어 주는 고리를 마련하려는 뜻에서 생겼다. 동시에 산술이론에 대한 논의가 본질적으로 빠져 있는 브릴--뇌터의 보고\srcfn{1)}{\emph{Die Entwicklung der Theorie der algebraischen Funktionen in älterer und neuerer Zeit}, \emph{Jahresb. d. D. Math. Verein.} 3 (1894) [이하 브릴--뇌터].}를 보완하는 글로 볼 수 있다\srcfn{2)}{크로네커의 판별식 연구에 대한 논의는 예외이다(\emph{J. f. M.} 91).}.

브릴--뇌터 외에도 각 이론에 대해서는 다음 백과사전 논문들과 거기에 실린 문헌을 참고할 수 있다.

리만 이론: Wirtinger (II B 2), 1--19번;

바이어슈트라스 이론: Wirtinger (II B 2), 23, 32, 33, 34번;

브릴--뇌터의 대수기하학적 이론: Wirtinger (II B 2), 21--31번; Berzolari (III C 4), 23--38번;

데데킨트--베버 및 헨젤--란츠베르크의 산술이론: Hensel (II C 5) [이하 헨젤]\srcfn{3)}{데데킨트--베버에 관한 부분은 A. Ostrowski가 썼다.};

더 일반적으로 대수적 원소의 산술이론: Landsberg (I C 5), 여러 변수 대수함수의 산술이론: Jung (II C 6).

수체론에 대해서는 힐베르트의 \glqq{}수론 보고\grqq{}\srcfn{4)}{\emph{Die Theorie der algebraischen Zahlkörper}, \emph{Jahresb. d. D. Math. Verein.} 4 (1894/95).}와 디리클레--데데킨트의 수론 강의 [이하 데데킨트--수론]\srcfn{5)}{Braunschweig, Vieweg, 제4판, 1894.}를 참조하라.

\begin{center}
\textbf{목차}
\end{center}
\begin{enumerate}
\item 여러 이론의 기초.
\item 대수적 수와 대수함수의 평행성(켤레원소, 기저정리, 모듈 개념).
\item 대수적 수와 대수함수의 아이디얼론에서 나타나는 평행성과 차이.
\item 급수전개의 산술적 기초와 아이디얼론의 관계.
\item 대수함수의 산술이론과 기하학적 이론의 비교. 서로 다른 불변성 개념.
\item 비교의 계속. 특이점.
\item 잔여정리와 아이디얼론 정리들의 비교.
\item 대수함수와 대수적 수에서의 초월적 문제.
\end{enumerate}

\begin{center}
\textbf{1. 여러 이론의 기초}
\end{center}

리만의 이론은 \emph{순전히 초월적}이며 \emph{존재정리}(디리클레 원리, 슈바르츠--노이만 방법)에 바탕을 둔 유일한 이론이다. 여기서는 대수함수의 \emph{적분}이 일차적이고 대수함수 자체가 이차적이다. 리만곡면 위의 함수는 영점과 극점으로 특징지어지는데, 이는 산술이론의 폴리곤(또는 인자)에 대응한다. 한편 바이어슈트라스도 함수 자체를 그 영점 및 극점과 함께 다룬다. 존재정리 외의 주된 문제는 주어진 극점들을 갖는 모든 함수를 세우는 것이며, 이는 제1종 및 제2종 적분을 통해 해결된다. 더 나아가 극점들을 지정한 뒤에도 들어가는 임의상수의 수를 묻는다. 이 마지막 문제에는 리만--로흐 정리가 답한다. 산술이론에서는 폴리곤류(또는 인자류)의 차원을 묻는 문제가, 대수기하학적 이론에서는 완전 선형계의 차원을 묻는 문제가 이에 대응하며, 여기서도 리만--로흐 정리가 답을 준다. 함수의 실제 구성은 산술이론에서는 기저정리로, 대수기하학적 이론에서는 잔여정리로 이루어진다. 특히 제1종 적분의 피적분함수의 영점들에는 산술이론의 미분류가 대응하고, 대수기하학적 이론의 특수점군 이론, 즉 \((n-3)\)차 수반곡선인 \(\varphi\)-곡선이 잘라 내는 점군들의 선형계가 대응한다.

리만 이론은 역사적으로 다른 이론들보다 앞섰지만 존재정리들의 엄밀한 증명은 뒤에 이루어졌다. 마찬가지로 훗날에도 그 초월적 도구로, 순수 대수적 또는 산술적 방법으로는 오늘날까지 접근할 수 없는 대수적 문제들을 해결하였다. 여기서는 무엇보다 특이대응에 관한 후르비츠의 이론을 들어야 한다\srcfn{1)}{A. Hurwitz, \emph{Über algebraische Korrespondenzen und das erweiterte Korrespondenzprinzip}, \emph{Ann.} 28. 브릴--뇌터 제10절 참조.}.

대수적 수체론의 초월적 문제들과의 유추는 마지막 절에서 논의한다.

나머지 이론들에서는 \emph{대수함수}가 일차적이고 적분이 이차적이다. \emph{바이어슈트라스 이론}은 거듭제곱급수 전개와 유수정리라는 초월적 기초 위에 서지만, 그 뒤에는 고려되는 모든 함수를 명시적으로 제시하면서 대수적으로 전개된다.

헨젤--란츠베르크의 형태를 취한 \emph{산술이론}도 \emph{급수전개라는 초월적 기초}를 쓴다. 그러나 이 기초 위에서 각 점에 인자를 대응시킨 뒤에는 순수 산술적으로 진행하므로, 바이어슈트라스와 데데킨트--베버의 융합으로 볼 수 있다. 후자의 이론과 달리, 고려되는 모든 기저함수를 실제로 구성하는 방법도 함께 준다. 세부적으로는 특히 \glqq{}분수\grqq{} 인자를 도입함으로써 데데킨트--베버보다 여러 점에서 단순해진다. 최근 헨젤\srcfn{1)}{\emph{Mathematische Zeitschrift} 4; 그의 보고도 이 기초에 바탕을 둔다.}은 수렴증명을 위한 우세급수법만 제외하고 급수전개 이론에 산술적 기초를 주었으며, 이로써 이 이론은 거의 순수 산술적이 되었다.

브릴--뇌터의 \emph{대수기하학적 이론}은 한 가지의 \glqq{}연속하는 점들\grqq{}이라는 개념만을 전제로 한다. 이는 \emph{보통점}에서의 급수전개로 환원되고\srcfn{2)}{연속하는 \glqq{}특이\grqq{}점에 대해서는 브릴--뇌터 제6절, 특히 10, 11번을 참조하라. 보통의 연속점은 이를테면 접선과 곡선의 교점 또는 고차 접촉점(분기점)이다.} 바이어슈트라스의 원소 개념에 대응한다. 그 뒤에는 대수적 유리방법, 본질적으로 삼원 소거론의 방법을 써서 진행하며, 이 이론도 명시적 표현에까지 이른다.

데데킨트--베버의 원래 \emph{순수 산술이론}\srcfn{3)}{\emph{J. f. M.} 92 [이하 D.--W.].}은 그 정리들이 다시 존재증명의 성격을 띤다는 점에서 리만 이론과 가깝고, 방법의 완전한 순수성을 지킨다\srcfn{4)}{그 때문에 이후의 비교에서는 가끔만 나타나는 헨젤--란츠베르크 이론보다 다른 이론들과 비교하기에 더 알맞다. 다만 아벨 적분과 그 역에 관한 이론은 연속성을 전제하는 한 D.--W.에서도, 대수이론의 후속 발전(M. Noether, \emph{Ann.} 37)에서도 다루어지지 않았다.}. 변수들이 단지 부정원의 역할만 하는 이 산술적 기초는 제2절에서 더 자세히 밝히듯 \emph{대수적 수}의 이론과 평행하게 진행된다\srcfn{1)}{헨젤에서는 헨젤--란츠베르크 이론과 수론, 특히 \(p\)-진수론의 비교를 볼 수 있다.}. 이 기초는 점 개념을 산술적으로 도입할 수 있게 하며(헨젤 10a 참조), 그 결과 연속성에 대한 어떤 고찰도 없이 절대 리만곡면의 개념을 세울 수 있게 한다. 이로써 이론은 형식적인 성격을 얻는다. 예를 들어 미분의 개념도 순전히 형식적으로 도입된다. 점 개념으로 얻은 기초 위에서는, 원소 개념을 제외하면 역시 초월적 도구를 쓰지 않는 대수기하학적 이론의 방법 및 개념과 가까운 친족관계를 밝힐 수 있다. 이를 뒤의 절들에서 수행하겠다.

\begin{center}
\textbf{2. 대수적 수와 대수함수의 평행성\newline
(켤레원소, 기저정리, 모듈 개념)}
\end{center}

대수적 수체는 유리수체 \(R\) 위의 \(n\)차 체로 정의되고, 한 변수 대수함수체는 임의의 복소계수를 갖는 한 부정원 \(z\)의 모든 유리함수체 \(K(z)\) 위의 \(n\)차 체로 정의된다. 따라서 기초체의 특수한 성질을 완전히 도외시하는 모든 정리, 즉 기저, 노름, 대각합, 판별식에 관한 정리들은 두 체에 공통으로 성립한다\srcfn{2)}{D.--W., \S{}1, 2; 헨젤 4번(주 5 포함).}. 데데킨트--베버에서는 이 개념들을 모두 켤레원소 없이 도입한다. 그러나 슈타이니츠가 크로네커를 이어 보였듯, \emph{켤레체의 개념은 완전히 형식적으로 정의할 수 있다}.\srcfn{3)}{Steinitz, \emph{Algebraische Theorie der Körper}. (\emph{J. f. M.} 137, \S{}6, 8.)} 실제로 \(K\)가 임의의 체이고 \(\varphi(x)\)가 \(K\)에서 기약인 \(x\)의 다항식이면, \(\varphi(x)\)를 법으로 한 잉여류계는 체를 이룬다. \(x\)의 잉여류에 원소 \(j\)를 대응시키면, 동형대응 아래 다항식 \(f(x)\)가 대표하는 잉여류에는 원소 \(f(j)\)가 대응한다. \(\varphi(x)\)로 나누어지는 모든 다항식은 영류를 대표하므로 \(\varphi(j)\)는 영류에 대응한다. 따라서 \(j\)를 기호적으로 \(\varphi(x)=0\)의 근으로 볼 수 있고, 이것이 대수체 \(K(j)\)를 생성한다. 이 절차를 반복하면 \(n\)개의 켤레체 \(K(j),K(j'),\ldots,K(j^{(n-1)})\)를 얻는다. 그러므로 특히 대수함수에서도 \emph{순수 산술이론을 벗어나지 않고} 켤레원소를 말할 수 있다.

두 기초체 \(R\)과 \(K(z)\)에는 모든 \emph{정수적} 원소의 계가 정의되어 있다는 공통점이 있다. 각각 유리정수와 한 부정원의 다항함수이다. 이 정수적 원소들은 임의의 두 원소 \(a,b\)가 \emph{최대공약원} \(d\)를 가지며
\[
d=ma+nb
\]
로 쓸 수 있다는 특징을 갖는다. 여기서 \(d,m,n\)도 각각 \(R\) 또는 \(K(z)\)의 정수적 원소이고, \(d\)는 이 형태로 나타낼 수 있는 절댓값이 가장 작은 수 또는 차수가 가장 낮은 다항식이다. 대수체 또는 상체의 정수적 원소에 관한 정리들은 오직 이 성질에만 바탕을 두며, 이 성질은 유리정수 또는 한 부정원의 다항식에 대한 \emph{소인수분해의 유일성}을 낳는다. 최대공약원의 존재를 보이는 논증은 대수체의 정수적 원소들에 대한 \emph{기본계}, 즉 \emph{기저}의 존재도 준다. 실제로 \emph{정수적} 원소들로 이루어진 체의 기저들만 생각하면 그 판별식은 유리정수 또는 한 부정원의 다항식이 된다. 절댓값이 가장 작은 수 또는 차수가 가장 낮은 함수를 주는 모든 기저가 기본계를 이룬다.

더 일반적인 체까지 포함하여 기저 문제를 더 정확히 이해하려면 \emph{모듈 개념}을 도입한다. 문헌에서 원소의 성질에 따라 서로 다르게 나타나는 정의들을 포괄하도록 다음처럼 정하겠다. \emph{기초정역 \(\mathfrak S\)에 관한 모듈}이란, \emph{두 원소의 차가 다시 그 계에 속하고 한 원소와 \(\mathfrak S\)의 임의의 원소의 곱도 그 계에 속하는 원소계}를 말한다. \(\mathfrak S\)가 유리정수 전체이면 둘째 조건은 첫째 조건에서 따라온다\srcfn{1)}{모듈 개념은 데데킨트에게서 비롯되었다. 그는 『수론』 제2판에서 수모듈(\(\mathfrak S=\) 유리정수)을 처음 도입했으며, 정수계수 선형형식의 모듈도 거기에 암묵적으로 포함되어 있다. 데데킨트--베버에서는 원소가 대수함수이고 \(\mathfrak S\)가 한 부정원의 모든 다항식으로 이루어진 \glqq{}함수모듈\grqq{}이 나타난다. \(n\)개 부정원의 동차형식, 더 일반적으로 다항식의 모듈은 힐베르트(\emph{Annalen} 36)에게서 처음 나타난다. 여기서는 \(\mathfrak S\)가 모든 다항식으로 이루어지고 원소들 역시 다항식이므로 이 모듈은 우리의 아이디얼 정의에 들어간다. 크로네커의 다항식 \glqq{}모듈계\grqq{}는 추상적 정의가 아니라 기저에서 출발한다.}.

유한개의 원소로 이루어진 기저를 가져서 모든 원소를 \(\mathfrak S\)의 계수로 선형표현할 수 있는 모듈을 \emph{유한}(유한생성) 모듈이라 한다. 특히 일생성 모듈은 \(c\alpha\) 꼴이며, 여기서 \(\alpha\)는 기저원소이고 \(c\)는 \(\mathfrak S\)의 모든 원소를 돈다. \emph{원소 자체가 \(\mathfrak S\)에 속하는 모듈을 \(\mathfrak S\)의 아이디얼이라 한다.} 여기서 유한 아이디얼, 특히 일생성 아이디얼인 주아이디얼의 개념이 따라온다.

정수적 원소와 아이디얼에 관한 모든 기저정리는 다음 정리에 바탕을 둔다.

\emph{\(M\)이 기초정역 \(\mathfrak S\)에 관한 모듈이고, 그 원소들이 \(\mathfrak S\)의 계수를 갖는 \(n\)개 부정원의 선형형식들이라고 하자. \(\mathfrak S\)의 모든 아이디얼이 유한이면 \(M\)도 유한모듈이다. 특히 \(\mathfrak S\)의 모든 아이디얼이 주아이디얼이면 \(M\)은 선형독립인 형식들로 이루어진 기저를 갖는다.}\srcfn{1)}{이 정리는 여러 계수정역 \(\mathfrak S\)에 대하여 문헌에 따로 나타난다. 유리정수에 대해서는 힐베르트, \emph{Zahlbericht} \S{}3, 4를, 여러 부정원의 다항식에 대해서는 힐베르트, \emph{Über die vollen Invariantensysteme}, \S{}2 끝(\emph{Math. Ann.} 42)을 참조하라.}

실제로 \(M\)의 원소들이
\[
l(x)=a_1x_1+\cdots+a_nx_n
\]
으로 주어지면, 계수 \(a_n\) 전체는 \(\mathfrak S\)의 한 아이디얼을 돈다. 가정에 따라 이는 \(b_1a_n^{(1)}+\cdots+b_\rho a_n^{(\rho)}\) 꼴이다. 따라서 \(M\)에 속하는 선형형식
\[
m(x)=l(x)-b_1l^{(1)}(x)-\cdots-b_\rho l^{(\rho)}(x)
\]
은 \((n-1)\)개의 부정원에만 의존한다. 이를 유한번 반복하면 명제가 따라온다. 매 단계에서 \(\rho=1\)이면 기저는 각각 \(n,n-1,\ldots,1\)개의 부정원에 의존하는 \(n\)개의 형식으로 이루어지고, 따라서 영이 아닌 형식들은 선형독립이다.

이로부터 여러 경우의 기본계 정리를 다음과 같이 얻는다.

대수체가 정수적 대수원소 \(\delta\)를 기초체에 첨가하여 생기고, \(\mathfrak S\)가 기초체의 모든 정수적 원소로 이루어진다고 하자. 그러면 상체의 모든 정수적 원소는 다음 표현을 갖는다(예: \emph{Zahlbericht} \S{}3).
\[
\omega=\frac{a_1\delta^{n-1}+a_2\delta^{n-2}+\cdots+a_n}{\Delta(\delta)},
\]
여기서 \(a_i\)들은 \(\mathfrak S\)의 원소이고 \(\Delta(\delta)\)는 \(\delta\)의 판별식이다. 치환
\[
x_i=\frac{\delta^{n-i}}{\Delta(\delta)}
\]
에 의해 모든 정수적 원소의 모듈은 선형형식의 모듈로 바뀐다. 따라서 \(\mathfrak S\)의 모든 아이디얼이 유한인 모든 경우에 기본계가 존재한다. 같은 논증은 더 일반적으로 상체의 정수적 원소 정역에 있는 모든 아이디얼이 유한임을 보여 준다. 특히 유리정수 또는 한 부정원의 다항식에서는 모든 아이디얼이 주아이디얼이다. 그러므로 대수적 수체 또는 한 변수 대수함수체의 차수가 \(n\)이면 그 기본계는 \emph{정확히 \(n\)개의 원소}로 이루어진다. 또한 앞에서 본 바와 같이 이 체들의 모든 아이디얼이 유한이므로 상대체에도 기본계가 존재하며, 일반적으로 그 원소 수는 상대차수보다 많다\srcfn{1)}{\(\mathfrak S\)가 한 유한 대수적 수체의 정수환일 때 선형형식의 모듈에 관한 슈타이니츠의 연구에 따르면, 상대차수가 \(\nu\)일 때 \(\nu+1\)개의 기저원소면 충분하다(\emph{Math. Ann.} 71, \S{}4).}. 더 나아가 \(\mathfrak S\)가 여러 부정원의 모든 다항식으로 이루어진 정역일 때도 모든 아이디얼은 유한이므로\srcfn{2)}{힐베르트의 모듈기저정리(\emph{Math. Ann.} 36)에 따른다. 여기서 일관성을 위해 아이디얼이라 부른 다항식 모듈은 보통 \glqq{}형식모듈\grqq{} 또는 단순히 \glqq{}모듈\grqq{}이라 한다. 그러한 모듈 \(N\)이 유한이라는 힐베르트의 정리도 선형형식 정리로 환원할 수 있다. 실제로 \(f\)가 \(N\)에 속하는 \(\rho+1\)개 부정원 \(\xi_1,\ldots,\xi_\rho,\eta\)의 \(n\)차 다항식이고 항 \(\eta^n\)을 포함한다고 하자. 이는 선형변환으로 언제나 달성할 수 있다. 그러면 \(N\)의 모든 다항식은 \(f\)를 법으로 하여 \(a_1(\xi)\eta^{n-1}+\cdots+a_n(\xi)\) 꼴인 것과 합동이다. 따라서 \(\eta^{n-i}=x_i\)로 두면 선형형식의 모듈을 이룬다. 한 부정원에 대해서는 명제가 성립하므로 귀납적으로 \(\rho\)개 부정원의 정역에서 모든 아이디얼을 유한이라고 가정할 수 있다.}, 여러 변수 대수함수체에도 기본계가 존재한다. 이 경우에도 그 원소 수는 차수보다 많다.

\begin{center}
\textbf{3. 대수적 수와 대수함수의 아이디얼론에서 나타나는 평행성과 차이}
\end{center}

모든 아이디얼이 소아이디얼들의 거듭제곱곱으로 유일하게 표현된다는 아이디얼론의 기본정리는 대수적 수와 대수함수에 대하여 함께 증명할 수 있다. 유리정수 또는 한 부정원의 다항식에서는 모든 아이디얼이 주아이디얼이라는 성질만을 쓰면 된다\srcfn{3)}{Hurwitz, \emph{Über die Theorie der Ideale}, \emph{Gött. Nachr.} 1894의 서론에 있는 논평을 참조하라. 베버의 대수학 교과서에서는 크로네커를 따라 함수적을 도입하여 두 체에 평행한 증명을 준다. 크로네커 이론과 아이디얼론 사이의 연결은 \glqq{}확장된 가우스 정리\grqq{}가 담당한다(Hurwitz, 앞의 글).}. 증명의 핵심은 아이디얼 \(\mathfrak c\)가 아이디얼 \(\mathfrak a\)로 나누어진다는 것, 즉 \(\mathfrak c\)의 모든 원소가 \(\mathfrak a\)에 들어 있다는 것으로부터 곱표현 \(\mathfrak c=\mathfrak a\mathfrak b\)가 따라옴을 보이는 데 있다\srcfn{4)}{데데킨트의 『수론』, 보유 XI, 정리 VII, \S{}177에서 특히 강조되어 있다. 554쪽의 주도 참조하라. 이 정리는 대수체를 다룬다는 사실을 쓰며, 다항식의 아이디얼(모듈)에는 더 이상 성립하지 않는다. 그 경우 곱 대신 최소공배수가 나타난다.}. 후르비츠의 증명은 \glqq{}확장된 가우스 정리\grqq{}에 바탕을 둔다. 이 정리에 따르면 두 다항식 \(\varphi,\psi\)의 곱의 모든 계수를 나누는 정수적 원소 \(\omega\)는 \(\varphi\)와 \(\psi\)의 모든 계수의 곱도 나눈다. 데데킨트의 증명은 이 정리와 동치인 모듈정리에 바탕을 둔다\srcfn{1)}{『수론』, 보유 XI, 정리 VI, \S{}173. 두 정리의 동치성은 데데킨트가 지적하였다(\emph{Über die Begründung der Idealtheorie}, \emph{Göttinger Nachr.} 1895).}.

더 나아가 여모듈의 개념(헨젤 12번)과 그로부터 분기아이디얼(기초아이디얼, 디퍼런트) \(\mathfrak D\)을 함께 정의할 수 있다. 그 노름은 체판별식 \(D\)가 된다. 또한 다음 기본정리가 성립한다. \emph{분기아이디얼 \(\mathfrak D\)는 모든 주아이디얼 \(F'(\delta)\)의 최대공약아이디얼이다.} 여기서 \(\delta\)는 체의 모든 정수적 원소를 돌고, \(F(\delta)=0\)은 그때마다 대응하는 기약방정식이다. 여기서 \emph{체판별식은 소아이디얼의 제곱으로 나누어지는 유리소수 또는 \(z\)의 일차인자만을, 그리고 그 모두를 포함한다}는 것이 따라온다. \(F'(\delta)\) 자체에 대해서는
\[
F'(\delta)=\mathfrak f\mathfrak D,\qquad
\Delta(\delta)=R^2D
\]
라는 표현을 얻는다. 여기서 \(\mathfrak f\)는 \(\delta\)로부터 생성되는 환(오더, 정역)의 도체이다. 즉 \(\mathfrak f\)는 임의의 정수적 원소와 곱했을 때 모듈 \([1,\delta,\ldots,\delta^{n-1}]\)로 나누어지는 모든 정수적 원소로 이루어진다. \(R^2=\operatorname{Norm}\mathfrak f\)는 기저 \(1,\delta,\ldots,\delta^{n-1}\)을 정수적 원소들의 기저로 옮기는 치환행렬식의 제곱이다\srcfn{2)}{Dedekind, \emph{Über die Diskriminante endlicher Körper}, \emph{Abhandl. der Gött. Ges. d. Wissensch.} 29 (1882). 거기서 대수적 수에 대하여 준 정의들은 대수함수로 곧바로 옮길 수 있고, 증명들도 거의 그대로 옮길 수 있다. D.--W.의 증명 배열은 이와 다르다. 아래를 보라. 인용한 정리들에 대해서는 \emph{Zahlbericht}, \S{}31, 32도 참조하라. 거기서는 \(F'(\delta)\)를 정수 \(\delta\)의 \glqq{}디퍼런트\grqq{}라 부른다.}.

뒤의 전개에서 생기는 \emph{차이}는 기초체의 특수한 성질들 때문이다. 대수적 수의 경우에는 \emph{기초체의 원소들이 동시에 지수}라는 사실 때문에, \(p\)에서 \(e\)제곱으로 나타나는 소아이디얼 \(\mathfrak p\)의 \((e-1)\)제곱보다 높은 거듭제곱이 분기아이디얼을 나눌 수 있다. 이는 \(e\)가 \(p\)로 나누어질 때 일어난다. 그 결과 관성체와 분기체라는 더 나아간 개념들\srcfn{3)}{\emph{Zahlbericht}, \S{}39--45 참조.}이 생기지만, 대수함수 이론에는 그 유사물이 없다. 무엇보다 대수함수에서는 기초체에 \emph{상수가 무한히 많이 존재}하므로 아이디얼류의 수가 유한하다는 성질이 사라지고, 이에 따라 유수의 결정과 관련된 모든 문제도 사라진다. 다만 어떤 의미에서는 모든 초월적 소형식의 정역\srcfn{1)}{Klein, \emph{Zur Theorie der Abelschen Funktionen}, \emph{Ann.} 36 (1890).}을 유체의 초월적 유사물로 볼 수 있다. 여기서는 각 점이 하나의 형식, 즉 한 주아이디얼로 생성되기 때문이다.

반대로 기초체 \(K(z)\)에 \emph{상수가 무한히 많이 존재}한다는 사실은 대수함수에서 대수적 수에는 유사물이 없는 정리와 문제를 낳는다. 먼저 증명의 배열이 어느 정도 단순해진다. 예를 들어 D.--W.가 아이디얼의 소아이디얼 분해 유일성을 증명할 때, 각 일차형식 \((z-c)\)가 모든 상수라는 무한히 많은 잉여류를 갖는다는 사실을 이용하여 \glqq{}확장된 가우스 정리\grqq{}나 그와 동치인 정리를 피할 수 있다. 반면 소수 \(p\)를 법으로 한 잉여류의 수는 유한하다\srcfn{2)}{D.--W. \S{}9. 212, 213쪽의 주를 참조하라.}.

실제로 새로운 결과를 낳는 것은 다음 사실이다. 계수가 \(z\)의 다항식인 모든 다항식은 \((z-c)\)를 법으로 일차인자들로 분해된다. 정수계수 다항식을 \(p\)를 법으로 볼 때에는 성립하지 않는 성질이다. 이 사실은 모든 상수 대신 모든 대수적 수만을 취해도 여전히 성립한다\srcfn{3)}{이 경우에도 이론 전체가 유지된다는 점을 D.--W.가 서론에서 지적한다.}. 여기서 \emph{모든 소아이디얼은 일차 아이디얼}이라는 것이 따라오고\srcfn{4)}{D.--W., \S{}9, 7번.}, 더 나아가 \emph{체판별식 \(D\)는 모든 판별식 \(\Delta(\delta)\)의 최대공약원}이 된다. 여기서 \(\delta\)는 모든 정수적 원소를 돈다\srcfn{5)}{D.--W., 224쪽.}. 두 사실 모두 대수적 수의 경우와 대조된다. 그 뒤에는 다시 분기아이디얼이 모든 주아이디얼 \(F'(\delta)\)의 최대공약아이디얼임이 따라온다.

또 다른 차이는 대수적 수에서는 유리수체 \(R\)가 \emph{절대적으로 특별한} 기초체, 즉 그 체에 포함된 \emph{소체}(단위원으로 생성되는 체)인 반면, \emph{기초체 \(K(z)\)는 상대적}이라는 데서 생긴다. 상수가 아닌 임의의 함수 \(\xi\)를 독립\emph{변수}로 택할 수 있으며, 그러면 원래 체는 \(K(\xi)\) 위의 대수체가 된다. 모든 소아이디얼 \(\mathfrak p\)가 일차 아이디얼이라는 정리, 즉 모든 함수가 \(\mathfrak p\)를 법으로 한 상수와 합동이라는 정리를 임의의 함수를 독립변수로 택할 수 있다는 사실과 결합하면, 대수적 수의 경우를 넘어서는 가장 중요한 개념인 \emph{점 개념}을 순수 산술적 형태로 얻는다(헨젤 10a 참조). \emph{각 점은 함수체와 동형인 한 상수체로 표현된다}\srcfn{1)}{잘 알려져 있듯 두 체가 \glqq{}동형\grqq{}이라는 것은 합, 차, 곱, 몫도 대응하도록 일대일로 서로 대응시킬 수 있다는 뜻이다.}. 이렇게 정의한 모든 점이 \glqq{}절대 리만곡면\grqq{}을 이룬다. 이 점 정의는 체에만 의존하고 특정 변수에는 의존하지 않는다. 다만 존재증명에서는 어떤 함수 하나를 독립변수 \(z\)로 특별히 택한다. 각 함수에 \(\mathfrak p\)를 법으로 한 상수 잉여를 대응시키면 소아이디얼 \(\mathfrak p\)가 점 \(P\)를 생성한다. 특히 \(\mathfrak p\)로 나누어지는 함수들은 \(P\)에서 사라지고, 그 역도 성립한다. \(\mathfrak p\)가 \(z\)에 관한 모든 소아이디얼을 돌면 \(z\)가 유한인 모든 점을 얻는다. 나머지 점들은 독립변수 \(\xi=1/z\)를 도입하여 얻으며, 주아이디얼 \((\xi)\)를 나누는 소아이디얼들에 대응한다. 당연히 점 개념에 이어지는 폴리곤, 폴리곤류, 폴리곤류의 차원 등의 개념도 대수적 수에는 유사물이 없다.

독립변수를 자유롭게 택할 수 있다는 사실은 분해 \(F'(\vartheta)=\mathfrak f\mathfrak D\)에 또 다른 해석도 준다. \(\vartheta\)를 독립변수로 보면, \(\vartheta\)로 생성되는 환은 \(z,\vartheta\)에 대하여 대칭적으로 정의되었으므로 이에 대응하여
\[
\left(\frac{\partial F}{\partial z}\right)=\mathfrak f\mathfrak a
\]
를 얻는다. 그리고 \(\mathfrak f\)는 두 주아이디얼 \(\partial F/\partial z\)와 \(\partial F/\partial\vartheta\)의 최대공약아이디얼이 된다. 즉 \emph{도체 \(\mathfrak f\)는 동시에 \(\vartheta,z\)에 관한 \glqq{}이중점 아이디얼\grqq{}이다}\srcfn{2)}{D.--W., 263쪽.}.

\begin{center}
\textbf{4. 급수전개의 산술적 기초와 아이디얼론의 관계}
\end{center}

헨젤--란츠베르크에서는 급수전개라는 초월적 도구와 함수론에서 가져온 점 개념이 아이디얼론을 대신한다. 최근 헨젤이 대수적 수와 대수함수의 급수전개에 부여한 산술적 기초\srcfn{3)}{\emph{Eine neue Theorie der algebraischen Zahlen}, \emph{Math. Zeitschr.} 2 (1918), 그리고 \emph{Neue Begründung der arithmetischen Theorie der algebraischen Funktionen einer Variablen}, 같은 곳 4 (1919). 더 자세한 서술은 헨젤 44번 이하를 참조하라.}는 따라서 아이디얼론과 동등한 것으로 볼 수 있다. 실제로 이 산술적 기초는 각 소아이디얼마다 그 소아이디얼이 주아이디얼이 되는 한 초월적 확대체로 넘어가서 통상적인 분해정리들이 성립하게 하는 데 귀착된다. 이 고찰은 한 소수 \(p\), 또는 일차함수 \(z-a\)를 나누는 모든 소아이디얼에 대하여 수행하면 충분하다.

먼저 \(z-a\), 또는 \(p\)의 정수 거듭제곱에 따른 모든 급수전개로 이루어지고 음의 거듭제곱은 매번 유한개만 허용하는 초월적 확대체를 도입한다. 후자는 \(p\)-진수체이다. 이 체에서 대수적 수체 또는 함수체를 정의하는 기약방정식은 \(\rho\)개의 기약인자의 곱으로 분해되며, 이는 \(p\) 또는 \(z-a\)가 \(\rho\)개의 \glqq{}으뜸\grqq{} 아이디얼 \(\mathfrak q_1,\ldots,\mathfrak q_\rho\)의 곱으로 분해되는 데 대응한다. 각 으뜸아이디얼을 다시 소아이디얼의 거듭제곱
\[
\mathfrak q_i=\mathfrak p_i^{\,l_i}
\]
으로 나타내는 데에는, 헨젤의 방법에서 초월적 확대체 위의 또 다른 대수체를 도입하는 것이 대응한다. 대수함수의 경우 이 체는 순수방정식으로 정의할 수 있고, 대수적 수의 경우에는 더 일반적으로 대수적으로 풀 수 있는 방정식으로 정의할 수 있다. 대수함수에서의 단순화는 이 경우 관성체와 분기체가 나타나지 않는다는 사실에 기인한다(제3절 참조).

이 연구들에서 소원소 \(\pi\)에 관하여 나타나는 \glqq{}종결급수\grqq{}\srcfn{1)}{\emph{Neue Begründung \(\ldots\)}, 식 (11).}
\[
v=c_0+c_1\pi+\cdots+c_{\sigma-1}\pi^{\sigma-1}-v_\sigma\pi^\sigma,
\]
여기서 \(c_0,\ldots,c_{\sigma-1}\)는 상수이고 \(v_\sigma\)는 체의 정수적 원소인데, 이와 정확히 대응하는 식을 D.--W.에서도 볼 수 있다\srcfn{2)}{D.--W., 231쪽 및 \S{}15.}. 거기서는 아이디얼론으로부터 이를 유도하며, 이 식이 한 점에서 함수가 영이 되는 차수 또는 무한대가 되는 차수를 정의하는 기초를 준다.

\begin{center}
\textbf{5. 대수함수의 산술이론과 기하학적 이론의 비교\newline
서로 다른 불변성 개념}
\end{center}

\glqq{}불변\grqq{}이라는 것은 \emph{오직 체 자체로 특징지어지는} 모든 개념을 말한다. 산술이론에서는 일반적으로 정의부터 체의 \emph{모든} 원소에 관하여 진술되고 어떤 기저나 특별히 선택한 변수에 관하여 진술되지 않으므로, 이 불변성이 이미 정의에 들어 있다. 앞에서 제시한 점 개념이 그 예이다. 반면 다른 이론들은 특정한 변수나 기저에서 출발하여 그 뒤에 그 개념이 이 선택과 무관함을 보인다. 여기서 불변성은 \emph{쌍유리변환에 대한 불변성}이 된다.

실제로 한편으로 체가 \(K(z)\)에 대수원소 \(s\)를 첨가하여 생기고, 다른 한편으로 \(K(\xi)\)에 \(\sigma\)를 첨가하여 생긴다고 하자. \(f(z,s)=0\)과 \(F(\xi,\sigma)=0\)은 각각 \(s,\sigma\)의 기약방정식이다. 더 일반적으로는 \(K(\eta)\)에 유한개의 원소를 첨가하고 대응하는 기약방정식들을
\[
g_1(\eta,\tau_1)=0,\qquad
g_2(\eta,\tau_1,\tau_2)=0,\ldots
\]
로 둘 수도 있다. 정의에 따라 모든 원소는 \(z,s\)로도, \(\xi,\sigma\)로도, \(\eta,\tau_1,\tau_2,\ldots\)로도 유리적으로 표현된다. 특히 \(z,s\)는 \(\xi,\sigma\)로 유리적으로 표현되고 그 역도 성립한다. 마찬가지로 \(z,s\)와 \(\eta,\tau_1,\tau_2,\ldots\) 사이에도 양방향의 유리표현이 있다. 따라서 평면곡선 \(f(z,s)=0\)은 \glqq{}쌍유리변환\grqq{}으로 평면곡선 \(F(\xi,\sigma)=0\)으로 옮겨지며, 또는 변수 \(\eta,\tau_1,\tau_2,\ldots\)의 공간에서 위 방정식들이 정의하는 곡선으로 옮겨진다. 불변성은 특정 곡선에 대하여 진술한 정의가 그 곡선에서 쌍유리변환으로 생기는 모든 곡선, 즉 임의 차원의 공간에 있는 모든 곡선에 대하여 유지된다는 방식으로 나타난다. 이들은 리만의 의미에서 한 \glqq{}류\grqq{}를 이룬다.

여기서 \glqq{}점\grqq{}은 \(f(a,b)=0\)을 만족하는 서로 결부된 값의 계 \(z=a,s=b\)로 정의된다. 이 점은 쌍유리변환으로 일반적으로 다시 \(F(\xi,\sigma)=0\) 위의 한 결부된 값의 계 \(\alpha,\beta\)로 옮겨진다. 그러나 \(a,b\)가 \glqq{}특이\grqq{}점인 특별한 경우에는 \(F(\xi,\sigma)=0\) 위의 여러 \glqq{}점\grqq{}에 대응할 수 있다. 또한 \(f(z,s)=0\)의 한 \glqq{}특이\grqq{}점이 유한개의 단순점, 즉 서로 떨어져 있거나 이웃한 점들에 대응하는 곡선을 언제나 줄 수 있다. 이것이 특이점 해소이다\srcfn{1)}{예를 들어 브릴--뇌터 제6절을 참조하라. \glqq{}특이\grqq{}점에 관한 자세한 내용은 다음 절에 있다.}. 마찬가지로 점군은 다시 점군으로 옮겨진다. 적절한 변환에서 특이점에 대응하는 단순점들의 수만큼 그 특이점을 세면 점의 수는 불변이다. 모든 곡선은 \glqq{}보통 다중\grqq{}점들만 갖는 곡선으로 쌍유리변환할 수 있으므로\srcfnmark{1)}, 기하학적 이론은 그러한 특별한 곡선들에 대해서만 개념들을 전개한 뒤, 가장 일반적인 특이점들을 낳는 변환까지 포함한 모든 쌍유리변환에 대하여 불변임을 보인다.

특히 초타원인 경우를 제외하면 이 \glqq{}류\grqq{}에서 특별한 곡선으로 \((p-1)\)차원 공간의 이른바 \glqq{}\(\varphi\)들의 정규곡선\grqq{}을 도입할 수 있다. 이 곡선은 특이점을 갖지 않는다. 여기서 \(\varphi\)들은 동차화한 \(f(z,s)=0\)에 대한 \(p\)개의 \glqq{}수반 \((n-3)\)차 곡선\grqq{}이다. \(f(z,s)=0\)을 \(F(\xi,\sigma)=0\)으로 옮기는 쌍유리변환에는 \(\varphi\)들의 선형변환이 대응하므로, 여기서는 사영적 의미의 불변성을 다루게 된다. 체의 모든 함수를 \(\varphi\)들의 유리함수로 표현하는 것을 기하학적 이론에서는 \glqq{}불변표현\grqq{}이라 한다\srcfn{1)}{예를 들어 브릴--뇌터 제8절을 참조하라. 여러 변수 대수함수에서는 불변성을 선형변환에 대한 불변성으로 환원할 수 있는가라는 문제가 다루어지지 않았다.}.

따라서 \(\varphi\)들의 \emph{선형계}, 또는 이들이 잘라 내는 \glqq{}특수점군\grqq{}의 선형계는 모든 쌍유리변환 아래 자기 자신으로 옮겨지며, 오직 체로만 특징지어진다. 이는 산술이론의 미분류에 대응한다. 산술이론에서는 다시 체의 \emph{모든} 함수를 독립변수로 보았을 때의 성질들로 정의하기 때문에 불변성이 정의에 직접 나타난다.

\begin{center}
\textbf{6. 비교의 계속. 특이점}
\end{center}

점 개념을 서로 다르게 이해한다는 사실과 관련하여, 대수기하학적 이론에서 특별한 어려움을 일으키는 곡선의 \glqq{}특이점\grqq{}은 산술이론의 기초에는 전혀 나타나지 않는다. 따라서 기하학적 이론에서 이론을 세우는 데 필요한 \glqq{}수반함수\grqq{}도 산술이론의 구성에는 필요하지 않고, 더 높은 특수 부분에만 속한다. 반대로 기하학적 이론의 기초에는 분기점이 나타나지 않는다. 실제로 아론홀트 정규형으로 적분을 표현할 때에도 분기점은 사라진다(헨젤 27번, 식 80).

제5절에서 준 특이점의 정의는 다음과 같이 쓸 수 있다. \(f(z,s)=0\), 또는 더 높은 차원 공간의 한 곡선이 \(z=a,s=b\), 또는 \(z=a,s_i=b_i\)에서 \emph{특이점을 갖는다는 것은, 이 값의 계를 절대 리만곡면의 둘 이상의 점}(산술적 의미의 점)\emph{이 취하거나, 하나 이상의 점이 더 높은 차수로 취한다는 뜻이다.} 후자의 경우에는 \(z\)의 분기폴리곤에도 \(s\)의 분기폴리곤에도 속하는 한 점을 다룬다. 이는 기하학적 이론의 \glqq{}연속하는 점들\grqq{}에 대응하며, 그 점은 \(f(z,s)=0\), 또는 더 높은 차원 공간의 곡선의 고차 첨점이 된다.

가정에 따라 체의 임의의 함수 \(\eta\)는 \(z,s\)로 유리적으로 표현된다. 그러므로 동형성 때문에 \(z=a,s=b\)가 되는 절대 리만곡면의 각 점은, \(z,s\)로 나타낸 각 함수 \(\eta\)에 \(z=a,s=b\)에서의 값을 대응시킴으로써 생긴다. 그 점이 \emph{특이}하려면 \(z=a,s=b\)에서 여러 값을 취하거나 어떤 값의 계를 여러 번 취하는 함수 \(\eta\)가 적어도 하나 있어야 한다. \(z,s\)에 의한 유리표현 때문에 이는 그 표현의 \emph{분자와 분모가 \(z=a,s=b\)에서 동시에 사라질 때에만} 가능하다. 여기서 \emph{위 정의는 \(z=a,s=b\)에서 \(f,\partial f/\partial z,\partial f/\partial s\)가 모두 사라진다는 통상적인 정의와 동일하다}는 것이 따라온다. 실제로 이 경우
\[
\eta=\frac{\partial f/\partial z}{\partial f/\partial s}
\]
는 말한 종류의 함수로서 \(z=a,s=b\)에서 여러 값을 갖는다. 반대의 경우에는 어떤 함수가 \(0/0\) 꼴이 되더라도 오직 한 값만을 취한다.

산술적 점 개념에서는 어떤 한 함수에 서로 다른 값의 계가 대응하기만 해도 두 점을 서로 다르다고 보기 때문에 특이점이 나타날 수 없다. 산술적 점 개념의 기초가 되는 아이디얼론에서도, 체의 정수적 원소 \emph{전체}를 언제나 동시에 고려하기 때문에 특이점 개념이 나타나지 않는다. 실제로 \(\delta\)가 \(z\)에 관하여 정수적인 \(F(z,\delta)=0\)의 유한영역 특이점들은 모두, 둘째 정의와 제3절 끝에 따라 \glqq{}이중점 아이디얼\grqq{} \(\mathfrak f\)로 생성된다. 한편 기본정리에 따라 분기아이디얼 \(\mathfrak D\)는 모든 주아이디얼 \(F'(\delta)\)의 최대공약아이디얼이다. 따라서 각 소아이디얼 \(\mathfrak p\)에 대하여 대응하는 \(\mathfrak f\)가 \(\mathfrak p\)로 나누어지지 않도록 \(\delta\)를 택할 수 있다. 또 \(\mathfrak f\)는 \(F'(\delta)\)와 \(F'(z)\)의 최대공약아이디얼이므로, 두 주아이디얼 \(F'(\delta),F'(z)\) 가운데 적어도 하나는 \(\mathfrak p\)로 나누어지지 않는다. 그러므로 \emph{모든 \(F'(\delta)\)와 \(F'(z)\)가 동시에 사라지는 값의 계는 없다}. 따라서 값의 계 \(z=a,\delta_i=b_i\)는 절대 리만곡면의 \emph{한} 점에서만 취해진다. 즉 무한차원 공간의 곡선으로 해석할 수 있는 \emph{체의 모든 정수적 대수함수 전체는 유한영역에서 특이점을 갖지 않는다}. 아이디얼론에서는 유한값만을 다룬다.

모든 정수적 원소의 기저함수 \(\omega_1,\ldots,\omega_n\) 역시 유한영역에서 특이점을 갖지 않는다. 앞에서 본 바와 같이 모든 \emph{정수적} 함수와 동형으로 대응하는 상수계만으로도 산술적 의미의 각 점은 이미 \emph{유일하게} 정해진다. 따라서 임의의 \emph{정수적} 함수 \(s_i\)들로 정의되는 \glqq{}공간곡선\grqq{}이 \(s_i=\alpha_i\)에서 유한영역의 특이점을 갖는다면, \(s_i=\alpha_i\)에서 여러 값의 계를 취하는 \emph{정수적} 함수 \(\eta\)가 적어도 하나 있어야 한다. 이제 \(\omega_1,\ldots,\omega_n\)이 모든 정수적 원소의 기저이면
\[
\delta=a_1\omega_1+\cdots+a_n\omega_n
\]
이고 \(a_i\)들이 정수적 유리함수이므로, \(\omega\)의 각 값의 계에는 \(\delta\)의 값의 계 하나만이 대응한다. 따라서 \(\omega_1,\ldots,\omega_n\)이 정의하는 공간곡선은 유한영역에 특이점을 가질 수 없다. 또한 \emph{\(z\)가 유한인 절대 리만곡면의 각 점은 \(\omega\)들과 동형으로 대응하는 상수계 하나만으로 이미 유일하게 정의된다. 기저함수들은 바로 \(s_i=\alpha_i\)에서 여러 값을 갖게 되는 함수 \(\eta\)들이다.} 그러므로 예컨대 기저 \(1,\delta,\ldots,\delta^{n-1}\) 대신 정수적 원소들의 기저를 도입하면 특이점의 어려움을 피할 수 있다.

기하학적 이론에서는 절대 리만곡면의 모든 점을 유일하게 생성하는 문제, 더 일반적으로 주어진 점들에서 극을 갖는 모든 함수를 구성하는 문제를 수반형식과 그 몫을 통해 해결한다. 형식 \(\psi\)가 수반이라는 것은 기초곡선 \(f(z,s)=0\)의 각 \(i\)중점에서 \(\psi=0\)이 적어도 \((i-1)\)중점을 갖는다는 뜻이다. 고차 특이점은 먼저 보통 다중점들로 해소해야 한다. \emph{\(f(z,s)=0\)의 한 특이점에서 가장 일반적인 함수 \(\eta=\psi/\chi\)의 분자와 분모가 모두 사라져야 한다}는 사실은 위에서 보았다. 그것들이 수반이어야 함은 전역적으로 유한인 미분들을 고찰하면 드러난다. 거꾸로 수반조건들이 가장 일반적인 함수를 구성하는 데 충분하다는 것은 뒤에서 더 자세히 다룰 \glqq{}잔여정리\grqq{}가 보여 준다. 따라서 수반함수들은 산술이론의 기저함수들을 대신한다.

모든 특이점이 유한영역에 놓인다고 가정하자. 이는 변환으로 언제나 달성할 수 있다. 그러면 산술적으로 한 \emph{수반형식}에는 체에서 \emph{\glqq{}이중점 아이디얼\grqq{} \(\mathfrak f\)로 나누어지는 함수의 분자}가 대응한다. 이로부터 형식적으로 \emph{기저원소 \(\omega_1,\ldots,\omega_n\)을 수반형식의 몫으로 나타내는 표현}과, 그 결과 서로 다른 표현방식들의 관계를 다음처럼 얻는다.

\(R(z)\)가 기저 \(1,\delta,\ldots,\delta^{n-1}\)을 기저 \(\omega_1,\ldots,\omega_n\)으로 옮기는 치환행렬식이라 하자. 거꾸로 각 \(\omega\)는 \(z,\delta\)의 다항함수를 \(R(z)\)로 나눈 몫으로 표현된다.
\[
\omega_i=\frac{G_i(z,\delta)}{R(z)}.
\]
따라서 \(R(z)\omega_i\)는 \(\delta\)로 생성되는 환에 속하므로, 도체의 정의에 따라 \(R(z)\)는 \(\mathfrak f\)로 나누어져야 한다. \((R(z))^2=\operatorname{Norm}\mathfrak f\)에서는 \((R(z))^2\)에 대해서만 이 나눗셈이 따라온다. 한편 \(\omega_i\)는 정수적 함수이므로 유한한 모든 \(z\)에서 유한하게 남는다. 따라서 분자함수 \(G_i(z,\delta)\)도 \(\mathfrak f\)로 나누어져야 한다. 즉 \emph{모든 \(\omega_i\)의 분자와 분모는 수반이다.} \(\omega\)들의 이 표현으로부터 체의 모든 함수를 수반형식의 몫으로 나타내는 표현이 따라온다.

앞의 결과에 따르면 기하학적 이론은 \emph{한 원소 \(\delta\)로 생성되는 환}(오더)의 이론으로 특징지을 수 있다. 이는 \emph{모든 정수적 원소}를 동시에 고찰하는 산술이론과 대조된다. 따라서 환의 도체, 즉 특이점들이 결정적인 역할을 한다는 것이 분명하다. 환(오더)의 산술이론과 기하학적 이론 사이에는 더 많은 유추가 있다. 예를 들어 기하학적 이론은 수반곡선과의 교차에서 특이점으로 들어가는 교점들을 어디에서도 세지 않고, 그와 다른 점군들만을 고찰한다. 이에 대응하여 데데킨트\srcfn{1)}{\emph{Über die Anzahl der Idealklassen in den verschiedenen Ordnungen eines endlichen Körpers}, Braunschweig 기념논문 (1877), \S{}3--5. 이 절들에서 대수적 수에 관하여 진술한 정리들은 한 변수 대수함수에도 문자 그대로 성립한다.}는 도체아이디얼 \(\mathfrak f\)와 서로소인 환의 아이디얼들만을 고찰하고, 이 범위의 아이디얼에 대하여 소아이디얼 분해의 유일성을 포함한 모든 나눗셈법칙을 세운다.

기하학적 이론의 기초가 되는 \glqq{}대수함수의 기본정리\grqq{}\srcfn{2)}{M. Noether, \emph{Über einen Satz aus der Theorie der algebraischen Funktionen}, \emph{Math. Ann.} 6 (1873).}, 또는 적어도 그 직접적인 한 귀결도 어떤 의미에서는 환의 산술이론에서 유사물을 갖는다. 이 정리는
\[
\psi(z,s)=A(z,s)f(z,s)+B(z,s)\varphi(z,s)
\]
로 표현할 수 있기 위한 조건들을 제시한다. 여기서 나타나는 모든 양은 \(s,z\)의 다항식이며, 더 일반적으로 \(x_1,x_2,x_3\)의 정수적 동차형식이다. 이 조건들에서\srcfn{3)}{브릴--뇌터, 55번.}, \(f(z,s)=0\)을 이용할 때 몫 \(\psi/\varphi\)가 \(f\)에 수반이면 언제나 \(s,z\)의 다항함수 \(B(z,s)\)와 같아진다는 것이 따라온다. 이에 대응하는 것은 이 모든 비교의 기초가 되는 제3절의 정리, 즉 \(\delta\)로 생성되는 환의 도체 \(\mathfrak f\)가 \(f(z,\delta)=0\)의 이중점 아이디얼과 같다는 정리이다. 실제로 도체의 정의에 따라 \(\mathfrak f\)로 나누어지는 모든 정수적 대수원소는 환에 속하므로 \(z,\delta\)의 다항식으로 표현된다. 반면 방금 말한 정리는 그러한 원소가 수반함수라고 말한다\srcfn{4)}{여기서는 데데킨트를 따라 \(\delta\)가 정수적 대수원소라는 가정을 계속 두었다. 이는 변환으로 언제나 달성할 수 있다. 기하학에 더 정확히 대응하여 \(s\)가 \(z\)에 대하여 대수적 분수함수일 수도 있는 \(f(z,s)=0\)의 경우는 헨젤--란츠베르크가 다룬다. 헨젤 26, 29번 참조.}.

\begin{center}
\textbf{7. 잔여정리와 아이디얼론 정리들의 비교}
\end{center}

방금 말한 \glqq{}대수함수의 기본정리\grqq{}의 귀결은 기하학적 이론에서 곧바로 \glqq{}잔여정리\grqq{}로 이어지며, 그 결과 기하학적 이론 전체가 세워지는 기초로 이어진다. 그러나 어떤 면에서 잔여정리는 다른 기초들보다 더 초등적이다. 산술적으로는 소아이디얼 분해의 유일성, 더 정확히는 그 분해정리의 바탕에 있는 사실들의 직접적인 귀결로 표현할 수 있으므로 더 높은 환론에 의존하지 않기 때문이다. 다음과 같이 하면 이 산술적 표현에 도달한다.

변수들을 분수선형변환하여 기초곡선 \(f(z,\delta)=0\)의 특이점들과 고려하는 유한개의 점군이 모두 유한영역에 놓이게 하자. 또한 그 뒤에 변수들의 정수계수 선형변환을 가하여 \(\delta\)가 \(z\)에 관하여 정수적이게 하자. 그러면 절대 리만곡면에서 \(z=a,\delta=b\)가 되는 점 \(P\)를 생성하는 소아이디얼 \(\mathfrak p\)에는 \(f(z,\delta)=0\) 위의 점 \(z=a,\delta=b\)가 대응한다. 더 일반적으로 아이디얼
\[
\mathfrak a=\mathfrak p_1^{\rho_1}\cdots\mathfrak p_\sigma^{\rho_\sigma}
\]
은 각 점 \(z=a_i,\delta=b_i\)를 대응하는 중복도만큼 포함하는 점군을 생성한다. \(f(z,\delta)=0\) 위의 유한영역에 있는 모든 점이 이렇게 생성된다. 한 점 \(P\)에서는 \(\mathfrak p\)로 나누어지는 모든 함수 \(g_1,g_2,\ldots\)가 사라지고, 거꾸로 이 함수들 전체가 아이디얼 \(\mathfrak p\)를 생성한다. 따라서 \(f\) 위의 대응점은 \(g_1(z,\delta)=0,g_2(z,\delta)=0,\ldots\)과 \(f=0\)의 교점이다. 아이디얼 \(\mathfrak a\)에 대응하는 점군도 마찬가지이다. 모든 아이디얼이 두 주아이디얼의 최대공약아이디얼이라는 정리에 따라, 한 점군을 잘라 내는 데에는 언제나 그러한 곡선 두 개면 충분하다. 특히 한 주아이디얼에 대응하는 점군은 \emph{하나의} 곡선 \(g(z,\delta)=0\)이 잘라 내며 그 역도 성립한다. 주아이디얼은 완전교차에 대응한다.

기하학적 이론에서는 두 점군 \(A,B\)가 같은 점군 \(R\)과 합쳐져 완전교차를 이룰 때 이들을 \emph{공잔여}라고 부른다. 이 점군들을 각각 아이디얼 \(\mathfrak a,\mathfrak b,\mathfrak r\)가 생성한다면
\[
\mathfrak a\mathfrak r=(\alpha),\qquad
\mathfrak b\mathfrak r=(\beta)
\]
라는 관계가 성립하며, 이는 아이디얼 \(\mathfrak a,\mathfrak b\)가 동치임을 보여 준다. 거꾸로 두 아이디얼의 동치 \(\mathfrak a=\eta\mathfrak b\)로부터도, 모든 아이디얼은 다른 한 아이디얼을 곱하면 주아이디얼이 된다는 정리에 따라 언제나 이 관계가 따라온다\srcfn{1)}{데데킨트, 『수론』 \S{}181 참조. 실제로 \(\mathfrak a\mathfrak r=(\alpha)\)가 되도록 \(\mathfrak r\)을 택하면 \(\mathfrak b\mathfrak r=\eta^{-1}(\alpha)\)이다. \(\mathfrak b\mathfrak r\)은 정수적 원소들만을 포함하므로 \(\eta^{-1}(\alpha)\)도 그러하며, 따라서 이것도 주아이디얼 \((\beta)\)가 된다.}. 즉 \emph{동치인 아이디얼과 공잔여 점군은 동일한 개념이다.}

이제 잔여정리는 다음과 같이 말한다. \emph{\(A,B\)가 공잔여이면, 같은 잔여를 갖는 수반곡선들로 두 점군을 잘라 낼 수 있다. 단 특이점들은 제외한다.} 한 함수의 영점과 극점은 언제나 공잔여이므로, 이 정리는 동시에 체의 가장 일반적인 함수를 수반형식의 몫으로 표현할 수 있음을 증명한다. 아이디얼론의 언어로 쓰면, 잔여정리는 \(\mathfrak a,\mathfrak b\)뿐 아니라 \(\mathfrak a\mathfrak f,\mathfrak b\mathfrak f\)도 동치라는 사실에 불과하다. 여기서 \(\mathfrak f\)는 이중점 아이디얼이다. 실제로 위의 결과에 따라
\[
\mathfrak a\mathfrak f\mathfrak m=(\alpha),\qquad
\mathfrak b\mathfrak f\mathfrak m=(\beta)
\]
가 주아이디얼이 되게 하는 아이디얼 \(\mathfrak m\)이 언제나 존재한다. \(\alpha=0,\beta=0\)은 점군 \(A,B\), 더 정확히는 그 가운데 특이점과 다른 부분군 \(A',B'\)을 잘라 내는 수반곡선들이다\srcfn{2)}{\(\mathfrak a=\mathfrak f'\mathfrak a_1,\ \mathfrak b=\mathfrak f'\mathfrak b_1,\ \mathfrak f=\mathfrak f'\mathfrak f_1\)이고 \(\mathfrak a_1,\mathfrak b_1,\mathfrak f_1\)가 서로소라면, 이미 \(\mathfrak a_1\mathfrak f\mathfrak n,\mathfrak b_1\mathfrak f\mathfrak n\)이 주아이디얼이 된다. 여기서 \(\mathfrak n\)은 \(\mathfrak f\)와 서로소이다. 실제로 모든 아이디얼은 주어진 한 아이디얼과 서로소인 아이디얼을 곱하여 주아이디얼로 만들 수 있다(D.--W. 214쪽 또는 『수론』 \S{}178, 정리 XI).}. 따라서 잔여정리의 산술적 증명은 소아이디얼들에 점들을 대응시킬 수 있다는 사실과, 동치인 아이디얼들을 같은 하나의 아이디얼로 곱하여 주아이디얼들로 만들 수 있다는 사실에 있다.

공잔여 점군들은 같은 수의 점으로 이루어질 필요가 없다. 마찬가지로 동치인 아이디얼들도 같은 차수일 필요가 없다. 예를 들어 모든 주아이디얼은 서로 동치이고, 완전교차가 생성하는 모든 점군은 서로 공잔여이다. 동치인 폴리곤(인자)과의 이러한 차이\srcfn{3)}{폴리곤 \(A\)는 절대 리만곡면의 유한개의 점으로 이루어진다. 두 폴리곤 \(A,B\)가 한 함수 \(\eta\)의 영점과 극점이면 이들을 동치라고 한다. 그러면 \(\eta\)는 폴리곤의 몫 \(\eta=A/B\)으로 표현된다(D.--W. \S{}17, 18).}는 \(z\)가 무한대가 되는 점들이 \(z\)에 관한 소아이디얼로 생성되지 않는다는 사실에서 설명된다. 아이디얼의 몫
\[
\eta=\frac{\mathfrak a}{\mathfrak b}
\]
으로 함수를 표현하는 것은 \(\mathfrak a,\mathfrak b\)의 동치를 말하지만, 여기에는 \(z\)가 무한대가 되는 곳에 있는 \(\eta\)의 영점이나 극점이 나타나지 않는다. 예를 들어 주아이디얼 \((\alpha)\)는 \emph{정수적} 함수 \(\alpha\)의 \emph{영점}에 대응하는 소아이디얼들의 곱이다. \(\alpha\)의 극점들은 나타나지 않는다. \(\alpha\)는 \(z\)가 무한대가 되는 점에서만 무한대가 되기 때문이다. 이런 의미에서 아이디얼론은 기하학적 이론의 \emph{형식론}과 정확히 대응한다. 형식론에도 영점만 있고 극점은 없다.

기하학에서는 같은 차수의 형식들로 몫을 만들어 형식에서 함수로 넘어간다. 산술적으로 이에 대응하는 것은 아이디얼류에서 불변적으로 정의된 폴리곤류(인자류)로 넘어가는 일이다. 이제 더 이상 어떤 변수도 특별히 선택되지 않으므로, 함수를 폴리곤의 몫으로 나타내면 모든 영점과 극점이 드러난다. 아이디얼류, 즉 공잔여 점군의 류와 폴리곤류의 관계는 다음과 같다. 아이디얼 \(\mathfrak a,\mathfrak b\)가 생성하는 폴리곤을 \(\mathfrak A,\mathfrak B\)라 하고, \(z\)가 무한대가 되는 점들의 폴리곤, 즉 \(z\)와 \(z\)의 정수적 함수들의 분모폴리곤을 \(\mathfrak N\)이라 하자. \(\mathfrak a,\mathfrak b\)가 동치일 때 \(\mathfrak A,\mathfrak B\)가 동치인 것은 두 아이디얼의 차수가 같을 때 그리고 그때에만 성립한다. 일반적으로는 \(\mathfrak A\mathfrak N^\rho\)와 \(\mathfrak B\mathfrak N^\sigma\), \(\rho\ne\sigma\), 의 동치만이 따라온다. 따라서 아이디얼들, 또는 공잔여 점군들을 류로 나누는 것은 무한히 많은 폴리곤류를 한 류로 합치는 것이다. 이것은 폴리곤들을 \(\mathfrak N\)의 거듭제곱을 법으로 한 류들로 나누는 것으로도 볼 수 있다\srcfn{1)}{헨젤--란츠베르크, 제25강, 438쪽. 두 변수 \(z,s\)를 폴리곤의 몫으로 표현하면 각 변수를 이원 동차화하는 것으로 볼 수 있다. 특히 \(z,s\)가 같은 분모폴리곤을 갖도록 택하면 기하학에서 통상 쓰는 삼원 동차화가 생긴다.}.

서로 다른 수의 점으로 이루어진 공잔여 점군에 대한 잔여정리는 순수 기하학적 문제에서만 나타난다. 대수함수에 적용할 때에는 실제로 같은 수의 공잔여 점들, 즉 폴리곤류에 대응하는 경우만 나타난다. 한 점군이 생성하는 \glqq{}완전 선형계\grqq{}, 곧 같은 수의 점으로 이루어지고 그 점군과 공잔여인 모든 점군은 대응하는 폴리곤이 생성하는 폴리곤류와 정확히 일치한다. \glqq{}두 완전계의 합\grqq{}과 \glqq{}두 폴리곤류의 곱\grqq{}이라는 개념들도 일치한다. 그러한 완전계가 유한개의 선형독립인 점군만을 포함한다는 사실은 잔여정리의 직접적인 귀결이다. 잔여정리는 그 선형계의 차원, 즉 선형독립인 점군의 수를, 잔여군을 지나는 임의의 고정된 차수의 모든 수반곡선의 차원으로 환원한다. 이 차원을 실제로 정하는 리만--로흐 정리에 도달하기 위해 기하학적 이론은 잔여정리로부터 먼저 \glqq{}환원정리\grqq{}, 즉 고정점에 관한 정리\srcfn{1)}{브릴--뇌터, 58번.}를 이끌어 내고, 다시 거기서 리만--로흐 정리를 얻는다.

완전히 이에 대응하여 산술이론은 각 폴리곤류의 차원이 유한함을 보이고, 정수적 원소들의 특수한 기저인 정규기저를 통해 그 차원을 정한다. 이 기저는 여모듈을 이용하여 정의된다. 이로써 여모듈과 분기폴리곤의 관계에서 리만--로흐 정리를 얻는다. D.--W.의 서술에서는 환원정리도 나타나지만\srcfn{2)}{D.--W., \S{}30, 2번.}, 기하학적 이론에서처럼 기초로 쓰이지는 않는다. 따라서 제6절에서 언급한 기저와 수반함수의 대응 때문에, 여기에서도 두 이론의 보조수단들 사이에 어느 정도의 대응이 있음을 확인할 수 있다. 요약하면 두 이론의 개념 구성은 표현과 배열이 다르기는 해도 본질적으로 일치한다. 두 이론은 서로 독립적으로 생겼으며, 기하학적 이론이 약 10년 앞선다. 끝으로 리만--로흐 정리는 원래 리만과 로흐에게서, 또 바이어슈트라스에게서 계수정리로 나타났음을 덧붙인다. 여기서 다루는 것은 행렬
\[
(\varphi_i(x_j))
\]
의 계수이다. \(\varphi\)들은 \(p\)개의 수반 \((n-3)\)차 곡선을 돌고, \(x\)들은 잔여군의 점들을 돈다.

\begin{center}
\textbf{8. 대수함수와 대수적 수에서의 초월적 문제}
\end{center}

대수함수와 대수적 수의 초월적 문제들 사이의 유추는 특히 힐베르트가 지적하였다\srcfn{3)}{\emph{Mathematische Probleme}, 문제 12, \emph{Göttinger Nachrichten} 1900.}. 본질적으로 존재문제를 다룬다. 주어진 리만곡면, 따라서 주어진 분기점을 갖는 대수함수가 존재하는가라는 리만의 문제에는, 주어진 군과 판별식을 갖는 대수적 수체가 존재하는가라는 문제가 대응한다\srcfn{4)}{그로부터 자연스럽게 생기는, 군에서 출발하여 대수함수를 다루는 방법은 R. König이 더 일반적인 문제들의 특수한 경우로 수행하였다. 특히 \emph{Riemannsche Funktionen- und Differentialsysteme in der Ebene}, \emph{J. f. M.} 148, 그리고 \emph{Math. Ann.} 78, 79를 참조하라. König에서는 대수체 대신 \(K(z)\)에 관한 특수한 유한모듈이 나타나며, \glqq{}류의 성질\grqq{}만이 성립한다.}. 대수함수의 경우에는 존재증명을 일반적으로 수행할 수 있지만, 대수적 수체에서는 특수한 기초체, 즉 유리수체와 이차수체에 대해서, 그리고 아벨군에 대해서만 성공하였다\srcfn{1)}{삼차와 사차 방정식에서 나타나는 군들의 단순한 경우와 팔차 방정식에서의 몇몇 매우 특수한 군들만은 임의의 기초체에 대하여 다룰 수 있다. G. Bucht, \emph{Über einige algebraische Körper 8. Grades}, \emph{Arkiv f. Math., Astron. och Physik} 6, 30번. E. Noether, \emph{Gleichungen mit vorgeschriebener Gruppe}와 거기서 인용한 Seidelmann의 학위논문, \emph{Math. Ann.} 78도 참조하라.}. 여기서 문제되는 것은 유리수체 위의 모든 아벨수체가 단위근체들로부터 합성된다는 크로네커의 정리와, 상대 아벨수체에 대한 퓌터\srcfn{2)}{\emph{Math. Ann.} 75.} 및 헤케\srcfn{3)}{\emph{Math. Ann.} 76.}의 대응 정리들이다. 따라서 구하는 수체들은 지수함수와 모듈러함수라는 초월적 수단을 통해 실제로 주어진다. 이는 대수함수에서의 순수 존재증명과 대조된다.

대수함수의 존재증명에는 전역적으로 유한인 적분들의 존재가 선행한다. 이에 대한 유사물은 한 대수적 수체에서 소수 \(l\)에 관하여 정의되는 \glqq{}특이 일차수\grqq{}의 존재이다. 그 \(l\)제곱근들은 상대적으로 비분기이고 유체를 구성하는 최초의 벽돌이 된다\srcfn{4)}{Furtwängler, \emph{Reziprozitätsgesetze für Potenzreste mit Primzahlexponenten in algebraischen Zahlkörpern}, \emph{Math. Ann.} 67, 72.}. 이 존재증명에는 수체 안에 주어진 잉여지표를 갖는 소아이디얼이 언제나 존재한다는 정리가 본질적으로 쓰인다. 따라서 이 정리는 전역적으로 유한인 적분의 존재가 바탕을 두는 경곗값문제의 유사물로 볼 수 있다.

전역적으로 유한인 적분들은 두 폴리곤이 같은 류에 속하는지, 즉 같은 수의 점으로 이루어진 두 점군이 공잔여인지를 판정하는 초월적 기준을 준다. 이것은 \glqq{}공잔여 경로\grqq{}에 걸쳐 취한 제1종 적분들이 사라진다고 말하는 아벨 정리, 또는 그 역을 포함한 대응 미분정리이다\srcfn{5)}{예를 들어 M. Noether, \emph{Ann.} 37의 표현을 참조하라. 역사적으로 잔여정리가 아벨 정리에서 생겨났기 때문에 기하학적 이론은 적분의 합을 본떠 점군과 선형계의 \glqq{}합\grqq{}을 말한다. 이는 산술이론에서 류의 \glqq{}곱\grqq{}을 말하는 것과 대조된다.}.

대수적 수체에서 아벨 정리의 유사물은 \glqq{}\(l\)제곱잉여의 상호법칙\grqq{}이다. 이는 수체 자체에서, 또는 적어도 적절한 상체에서 두 아이디얼이 같은 류에 속한다는 판정기준을 준다. 그 내용은 이 상체에서 특이 일차수들이 같은 류의 모든 아이디얼에 관하여 같은 잉여지표를 갖는다는 식으로도 쓸 수 있다. 이 마지막 사실은 원래 체에서도 성립하지만, 그곳에서는 상호법칙과 일치하지 않는다\srcfn{1)}{Furtwängler, 앞의 글, \emph{Annalen} 67, 특히 서론 참조.}.

끝으로 헨젤의 \(p\)-진수론은 대수함수의 거듭제곱급수 전개에 대한 유사물을 준다. 이에 대해서는 헨젤의 보고를 참조하라.
